\documentclass[11pt,a4paper]{article}
\usepackage[utf8]{inputenc}
\usepackage[T1]{fontenc}
\usepackage[spanish,es-minimal,es-noshorthands]{babel}
\AtBeginDocument{\renewcommand{\sin}{\operatorname{sen}}\renewcommand{\arcsin}{\operatorname{arcsen}}}
\usepackage{lmodern}
\usepackage[sfdefault,lining]{FiraSans}
\usepackage[a4paper,top=2.3cm,bottom=2.2cm,left=1.8cm,right=1.8cm,headheight=1.3cm,headsep=0.4cm,footskip=1cm]{geometry}
\usepackage{amsmath,amssymb}
\usepackage{icomma}
\usepackage{graphicx}
\usepackage[table]{xcolor}
\usepackage{tikz}
\usepackage{fontawesome5}
\usepackage{ccicons}
\usepackage{fancyhdr}
\usepackage{lastpage}
\usepackage{tasks}
\usepackage{needspace}
\usepackage[most]{tcolorbox}
\usepackage{array}
\usepackage[hidelinks]{hyperref}
\AtBeginDocument{\hypersetup{pdftitle={Ficha de ejercicios de la unidad 8: Distribuciones bidimensionales. 1º Bach.},pdfsubject={Matemáticas CCSS 1º Bach.},pdfkeywords={pdfkeywords={distribuciones bidimensionales, tabla de doble entrada, nube de puntos, correlación, regresión, covarianza, coeficiente de correlación lineal, rectas de regresión}}}}

\definecolor{navy}{HTML}{1B2748}
\definecolor{accent}{HTML}{E8603A}
\definecolor{teal}{HTML}{1F9E90}
\definecolor{soft}{HTML}{EEF1F8}
\definecolor{muted}{HTML}{5D6685}

\setlength{\parindent}{0pt}
\setlength{\parskip}{3pt}
\settasks{label=\alph*), label-format=\bfseries\color{navy}, label-width=1.3em, item-indent=1.8em, column-sep=0.6em, after-item-skip=5pt, before-skip=3pt, after-skip=4pt}

% ---------- capçalera i peu ----------
\newcommand{\llicencia}{\href{https://creativecommons.org/licenses/by-nc-sa/4.0/deed.es}{\ccbyncsa}\ \ Licencia Creative Commons Atribución-NoComercial-CompartirIgual 4.0 Internacional (CC BY-NC-SA 4.0)}
\pagestyle{fancy}
\fancyhf{}
\fancyhead[L]{\raisebox{-0.35\height}{\includegraphics[height=0.95cm]{logo-gva.png}}}
\fancyhead[C]{\small\color{muted}Ficha de ejercicios · Unidad 8: Distribuciones bidimensionales · 1\textsuperscript{o} Bach.}
\fancyhead[R]{\raisebox{-0.35\height}{\includegraphics[height=0.95cm]{logo-vermellar.png}}}
\fancyfoot[L]{\scriptsize\color{muted}\llicencia}
\fancyfoot[R]{\small\color{muted}\thepage}
\renewcommand{\headrulewidth}{0.4pt}
\renewcommand{\headrule}{\color{muted!40}\hrule height \headrulewidth}
\fancypagestyle{primera}{\fancyhf{}\renewcommand{\headrulewidth}{0pt}
  \fancyfoot[L]{\scriptsize\color{muted}\llicencia}
  \fancyfoot[R]{\small\color{muted}\thepage}}

% ---------- nivells i exercicis ----------
\newcommand{\estrelles}[1]{{\color{accent}\ifcase#1\or\faStar\faStar[regular]\faStar[regular]\or\faStar\faStar\faStar[regular]\or\faStar\faStar\faStar\fi}}
\newcounter{exer}
\newcommand{\ex}[1]{\par\needspace{4\baselineskip}\medskip\refstepcounter{exer}\noindent{\color{navy}\bfseries\theexer.}\ \ }

\newcommand{\seccio}[2]{\par\needspace{16\baselineskip}\bigskip
  \begin{tcolorbox}[colback=navy,colframe=navy,coltext=white,arc=2mm,boxrule=0pt,left=3mm,right=3mm,top=1.2mm,bottom=1.2mm]
  {\large\bfseries #1\quad #2}
  \end{tcolorbox}\vspace{-2pt}}

\newcommand{\exf}[1]{\par\medskip\refstepcounter{exer}\noindent{\color{navy}\bfseries\theexer.}\ \ }
\providecommand{\gt}{>}\providecommand{\lt}{<}
% quadre de recordatori (només a la fitxa d'exercicis)
\newcommand{\recorda}[1]{\par\smallskip\begin{tcolorbox}[colback=accent!7,colframe=accent!55,arc=2mm,boxrule=0.5pt,left=2.5mm,right=2.5mm,top=1mm,bottom=1mm]\small{\color{accent}\faLightbulb[regular]}\ \textbf{Recuerda:}\ #1\end{tcolorbox}}
\newcommand{\res}[1]{#1}\newcommand{\rest}[1]{#1}

\def\fmneed{11}% que el títol d'una secció no quede sol al final del full
% =====================================================================
% Capçalera FentMates de les fitxes: un color per a cada curs.
% Els colors són els de la pàgina inicial (LEVELS a assets/js/materials.js);
% el logo ve de _eines/logo/pdf/ (python3 _eines/logo/genera.py).
%
% Ús, al final del preàmbul de la unitat:
%   % =====================================================================
% Capçalera FentMates de les fitxes: un color per a cada curs.
% Els colors són els de la pàgina inicial (LEVELS a assets/js/materials.js);
% el logo ve de _eines/logo/pdf/ (python3 _eines/logo/genera.py).
%
% Ús, al final del preàmbul de la unitat:
%   % =====================================================================
% Capçalera FentMates de les fitxes: un color per a cada curs.
% Els colors són els de la pàgina inicial (LEVELS a assets/js/materials.js);
% el logo ve de _eines/logo/pdf/ (python3 _eines/logo/genera.py).
%
% Ús, al final del preàmbul de la unitat:
%   \input{../comu/fentmates.tex}
%   \fmcurs{1eso}
% i, al començament del document:
%   \capcalera{Matemàtiques}{1r ESO}{Fitxa d'exercicis · Unitat 2}{Divisibilitat}
% Opcional abans de \input: \def\fmneed{14} (espai mínim abans de cada secció).
% =====================================================================
\makeatletter
\def\fm@logodir{../../logo/pdf}
\providecommand{\fmneed}{10}
\def\fm@cursnom{}

% clau del curs, color i color del text damunt del color
\newcommand{\fm@curs}[3]{\@namedef{fm@color@#1}{#2}\@namedef{fm@text@#1}{#3}}
\fm@curs{1eso}{E6B84A}{navy}
\fm@curs{2eso}{3B82F6}{white}
\fm@curs{3eso}{FF7A50}{white}
\fm@curs{4eso-a}{A78BFA}{white}
\fm@curs{4eso-b}{6D28D9}{white}
\fm@curs{1batx-mi}{9B1C1C}{white}
\fm@curs{1batx-ccss}{EF6B6B}{white}
\fm@curs{2batx-mii}{1E6B47}{white}
\fm@curs{2batx-ccss}{2BB3A3}{white}

\newcommand{\fmcurs}[1]{%
  \@ifundefined{fm@color@#1}{\PackageError{fentmates}{Curs desconegut: #1}{}}{}%
  \edef\fm@tmp{\noexpand\definecolor{curs}{HTML}{\@nameuse{fm@color@#1}}%
               \noexpand\colorlet{curstext}{\@nameuse{fm@text@#1}}}\fm@tmp
  \colorlet{cursfosc}{curs!72!black}%
  \def\fm@logo{\fm@logodir/vertical-#1.pdf}%
  \def\fm@marca{\fm@logodir/marca-#1.pdf}}

% ---------- capçalera de la primera pàgina ----------
% #1 matèria (petit), #2 curs, #3 tipus de document i unitat, #4 títol
\newcommand{\capcalera}[4]{%
  \gdef\fm@cursnom{#2}%
  \par\noindent
  \begin{tcolorbox}[colback=curs,colframe=curs,boxrule=0pt,arc=2.2mm,width=3.5cm,height=2.15cm,
      valign=center,left=3mm,right=2mm,top=1mm,bottom=1mm,nobeforeafter,box align=center]
    \color{curstext}{\footnotesize #1}\par\vspace{1pt}{\fontsize{19}{21}\selectfont\bfseries #2\par}%
  \end{tcolorbox}\hspace{5mm}%
  \begin{minipage}[c]{\dimexpr\linewidth-3.5cm-5mm-2.9cm\relax}
    {\color{cursfosc}\small\bfseries\MakeUppercase{#3}}\par\vspace{2pt}
    {\color{navy}\fontsize{20}{23}\selectfont\bfseries #4\par}
  \end{minipage}\hfill
  \begin{minipage}[c]{2.6cm}\raggedleft\includegraphics[width=2.6cm]{\fm@logo}\end{minipage}\par
  \vspace{2pt}}

% ---------- barres de secció del color del curs ----------
\renewcommand{\seccio}[2]{\par\needspace{\fmneed\baselineskip}\bigskip
  \begin{tcolorbox}[colback=curs,colframe=curs,coltext=curstext,arc=2mm,boxrule=0pt,left=3mm,right=3mm,top=1.2mm,bottom=1.2mm]
  {\large\bfseries #1\quad #2}
  \end{tcolorbox}\vspace{-2pt}}

% ---------- estrelles de dificultat del color del curs (a Batxillerat no se'n posen) ----------
\@ifundefined{estrelles}{}{%
  \renewcommand{\estrelles}[1]{{\color{curs}\ifcase#1\or\faStar\faStar[regular]\faStar[regular]\or\faStar\faStar\faStar[regular]\or\faStar\faStar\faStar\fi}}}

% ---------- capçalera i peu de les pàgines ----------
% primera pàgina: logos institucionals al peu
\fancypagestyle{primera}{\fancyhf{}\renewcommand{\headrulewidth}{0pt}
  \fancyfoot[L]{\raisebox{-0.32\height}{\includegraphics[height=0.78cm]{logo-gva.png}}\hspace{5mm}%
    \raisebox{-0.32\height}{\includegraphics[height=0.78cm]{logo-vermellar.png}}}
  \fancyfoot[R]{\scriptsize\color{muted}\href{https://creativecommons.org/licenses/by-nc-sa/4.0/deed.ca}{\ccbyncsa}\ \ CC BY-NC-SA 4.0 · fentmates.org\qquad\small\thepage}}
% la resta: el símbol a l'esquerra i el curs a la dreta, amb el títol al centre com abans
\fancyhead[L]{\raisebox{-0.3\height}{\includegraphics[height=0.62cm]{\fm@marca}}}
\fancyhead[R]{\small\bfseries\color{cursfosc}\fm@cursnom}
\renewcommand{\headrule}{\color{curs!45}\hrule height \headrulewidth}
\makeatother

%   \fmcurs{1eso}
% i, al començament del document:
%   \capcalera{Matemàtiques}{1r ESO}{Fitxa d'exercicis · Unitat 2}{Divisibilitat}
% Opcional abans de \input: \def\fmneed{14} (espai mínim abans de cada secció).
% =====================================================================
\makeatletter
\def\fm@logodir{../../logo/pdf}
\providecommand{\fmneed}{10}
\def\fm@cursnom{}

% clau del curs, color i color del text damunt del color
\newcommand{\fm@curs}[3]{\@namedef{fm@color@#1}{#2}\@namedef{fm@text@#1}{#3}}
\fm@curs{1eso}{E6B84A}{navy}
\fm@curs{2eso}{3B82F6}{white}
\fm@curs{3eso}{FF7A50}{white}
\fm@curs{4eso-a}{A78BFA}{white}
\fm@curs{4eso-b}{6D28D9}{white}
\fm@curs{1batx-mi}{9B1C1C}{white}
\fm@curs{1batx-ccss}{EF6B6B}{white}
\fm@curs{2batx-mii}{1E6B47}{white}
\fm@curs{2batx-ccss}{2BB3A3}{white}

\newcommand{\fmcurs}[1]{%
  \@ifundefined{fm@color@#1}{\PackageError{fentmates}{Curs desconegut: #1}{}}{}%
  \edef\fm@tmp{\noexpand\definecolor{curs}{HTML}{\@nameuse{fm@color@#1}}%
               \noexpand\colorlet{curstext}{\@nameuse{fm@text@#1}}}\fm@tmp
  \colorlet{cursfosc}{curs!72!black}%
  \def\fm@logo{\fm@logodir/vertical-#1.pdf}%
  \def\fm@marca{\fm@logodir/marca-#1.pdf}}

% ---------- capçalera de la primera pàgina ----------
% #1 matèria (petit), #2 curs, #3 tipus de document i unitat, #4 títol
\newcommand{\capcalera}[4]{%
  \gdef\fm@cursnom{#2}%
  \par\noindent
  \begin{tcolorbox}[colback=curs,colframe=curs,boxrule=0pt,arc=2.2mm,width=3.5cm,height=2.15cm,
      valign=center,left=3mm,right=2mm,top=1mm,bottom=1mm,nobeforeafter,box align=center]
    \color{curstext}{\footnotesize #1}\par\vspace{1pt}{\fontsize{19}{21}\selectfont\bfseries #2\par}%
  \end{tcolorbox}\hspace{5mm}%
  \begin{minipage}[c]{\dimexpr\linewidth-3.5cm-5mm-2.9cm\relax}
    {\color{cursfosc}\small\bfseries\MakeUppercase{#3}}\par\vspace{2pt}
    {\color{navy}\fontsize{20}{23}\selectfont\bfseries #4\par}
  \end{minipage}\hfill
  \begin{minipage}[c]{2.6cm}\raggedleft\includegraphics[width=2.6cm]{\fm@logo}\end{minipage}\par
  \vspace{2pt}}

% ---------- barres de secció del color del curs ----------
\renewcommand{\seccio}[2]{\par\needspace{\fmneed\baselineskip}\bigskip
  \begin{tcolorbox}[colback=curs,colframe=curs,coltext=curstext,arc=2mm,boxrule=0pt,left=3mm,right=3mm,top=1.2mm,bottom=1.2mm]
  {\large\bfseries #1\quad #2}
  \end{tcolorbox}\vspace{-2pt}}

% ---------- estrelles de dificultat del color del curs (a Batxillerat no se'n posen) ----------
\@ifundefined{estrelles}{}{%
  \renewcommand{\estrelles}[1]{{\color{curs}\ifcase#1\or\faStar\faStar[regular]\faStar[regular]\or\faStar\faStar\faStar[regular]\or\faStar\faStar\faStar\fi}}}

% ---------- capçalera i peu de les pàgines ----------
% primera pàgina: logos institucionals al peu
\fancypagestyle{primera}{\fancyhf{}\renewcommand{\headrulewidth}{0pt}
  \fancyfoot[L]{\raisebox{-0.32\height}{\includegraphics[height=0.78cm]{logo-gva.png}}\hspace{5mm}%
    \raisebox{-0.32\height}{\includegraphics[height=0.78cm]{logo-vermellar.png}}}
  \fancyfoot[R]{\scriptsize\color{muted}\href{https://creativecommons.org/licenses/by-nc-sa/4.0/deed.ca}{\ccbyncsa}\ \ CC BY-NC-SA 4.0 · fentmates.org\qquad\small\thepage}}
% la resta: el símbol a l'esquerra i el curs a la dreta, amb el títol al centre com abans
\fancyhead[L]{\raisebox{-0.3\height}{\includegraphics[height=0.62cm]{\fm@marca}}}
\fancyhead[R]{\small\bfseries\color{cursfosc}\fm@cursnom}
\renewcommand{\headrule}{\color{curs!45}\hrule height \headrulewidth}
\makeatother

%   \fmcurs{1eso}
% i, al començament del document:
%   \capcalera{Matemàtiques}{1r ESO}{Fitxa d'exercicis · Unitat 2}{Divisibilitat}
% Opcional abans de \input: \def\fmneed{14} (espai mínim abans de cada secció).
% =====================================================================
\makeatletter
\def\fm@logodir{../../logo/pdf}
\providecommand{\fmneed}{10}
\def\fm@cursnom{}

% clau del curs, color i color del text damunt del color
\newcommand{\fm@curs}[3]{\@namedef{fm@color@#1}{#2}\@namedef{fm@text@#1}{#3}}
\fm@curs{1eso}{E6B84A}{navy}
\fm@curs{2eso}{3B82F6}{white}
\fm@curs{3eso}{FF7A50}{white}
\fm@curs{4eso-a}{A78BFA}{white}
\fm@curs{4eso-b}{6D28D9}{white}
\fm@curs{1batx-mi}{9B1C1C}{white}
\fm@curs{1batx-ccss}{EF6B6B}{white}
\fm@curs{2batx-mii}{1E6B47}{white}
\fm@curs{2batx-ccss}{2BB3A3}{white}

\newcommand{\fmcurs}[1]{%
  \@ifundefined{fm@color@#1}{\PackageError{fentmates}{Curs desconegut: #1}{}}{}%
  \edef\fm@tmp{\noexpand\definecolor{curs}{HTML}{\@nameuse{fm@color@#1}}%
               \noexpand\colorlet{curstext}{\@nameuse{fm@text@#1}}}\fm@tmp
  \colorlet{cursfosc}{curs!72!black}%
  \def\fm@logo{\fm@logodir/vertical-#1.pdf}%
  \def\fm@marca{\fm@logodir/marca-#1.pdf}}

% ---------- capçalera de la primera pàgina ----------
% #1 matèria (petit), #2 curs, #3 tipus de document i unitat, #4 títol
\newcommand{\capcalera}[4]{%
  \gdef\fm@cursnom{#2}%
  \par\noindent
  \begin{tcolorbox}[colback=curs,colframe=curs,boxrule=0pt,arc=2.2mm,width=3.5cm,height=2.15cm,
      valign=center,left=3mm,right=2mm,top=1mm,bottom=1mm,nobeforeafter,box align=center]
    \color{curstext}{\footnotesize #1}\par\vspace{1pt}{\fontsize{19}{21}\selectfont\bfseries #2\par}%
  \end{tcolorbox}\hspace{5mm}%
  \begin{minipage}[c]{\dimexpr\linewidth-3.5cm-5mm-2.9cm\relax}
    {\color{cursfosc}\small\bfseries\MakeUppercase{#3}}\par\vspace{2pt}
    {\color{navy}\fontsize{20}{23}\selectfont\bfseries #4\par}
  \end{minipage}\hfill
  \begin{minipage}[c]{2.6cm}\raggedleft\includegraphics[width=2.6cm]{\fm@logo}\end{minipage}\par
  \vspace{2pt}}

% ---------- barres de secció del color del curs ----------
\renewcommand{\seccio}[2]{\par\needspace{\fmneed\baselineskip}\bigskip
  \begin{tcolorbox}[colback=curs,colframe=curs,coltext=curstext,arc=2mm,boxrule=0pt,left=3mm,right=3mm,top=1.2mm,bottom=1.2mm]
  {\large\bfseries #1\quad #2}
  \end{tcolorbox}\vspace{-2pt}}

% ---------- estrelles de dificultat del color del curs (a Batxillerat no se'n posen) ----------
\@ifundefined{estrelles}{}{%
  \renewcommand{\estrelles}[1]{{\color{curs}\ifcase#1\or\faStar\faStar[regular]\faStar[regular]\or\faStar\faStar\faStar[regular]\or\faStar\faStar\faStar\fi}}}

% ---------- capçalera i peu de les pàgines ----------
% primera pàgina: logos institucionals al peu
\fancypagestyle{primera}{\fancyhf{}\renewcommand{\headrulewidth}{0pt}
  \fancyfoot[L]{\raisebox{-0.32\height}{\includegraphics[height=0.78cm]{logo-gva.png}}\hspace{5mm}%
    \raisebox{-0.32\height}{\includegraphics[height=0.78cm]{logo-vermellar.png}}}
  \fancyfoot[R]{\scriptsize\color{muted}\href{https://creativecommons.org/licenses/by-nc-sa/4.0/deed.ca}{\ccbyncsa}\ \ CC BY-NC-SA 4.0 · fentmates.org\qquad\small\thepage}}
% la resta: el símbol a l'esquerra i el curs a la dreta, amb el títol al centre com abans
\fancyhead[L]{\raisebox{-0.3\height}{\includegraphics[height=0.62cm]{\fm@marca}}}
\fancyhead[R]{\small\bfseries\color{cursfosc}\fm@cursnom}
\renewcommand{\headrule}{\color{curs!45}\hrule height \headrulewidth}
\makeatother

\fmcurs{1batx-ccss}

\begin{document}
\thispagestyle{primera}

\capcalera{Matemáticas CCSS}{1\textsuperscript{o} Bach.}{Ficha de ejercicios · Unidad 8}{\raggedright Distribuciones bidimensionales}
\vspace{10pt}
Nombre y apellidos: \rule{6.9cm}{0.4pt}\hfill Grupo: \rule{1.6cm}{0.4pt}\hfill Fecha: \rule{2.2cm}{0.4pt}
\vspace{8pt}

\begin{tcolorbox}[colback=soft,colframe=soft,arc=2mm,left=3mm,right=3mm,top=1.5mm,bottom=1.5mm]
\small Resuelve los ejercicios en el cuaderno mostrando todos los pasos: para las variables bidimensionales, haz siempre la tabla con las columnas $x_i$, $y_i$, $x_i^2$, $y_i^2$ y $x_iy_i$. Redondea los resultados a dos o tres cifras decimales. Puedes comprobarlos con las calculadoras de los apuntes interactivos.
\end{tcolorbox}

% =====================================================
\seccio{1}{Distribuciones bidimensionales. Tablas de frecuencias}

\exf{1} Se ha preguntado a 25 familias el número de hijos ($X$) y el número de coches ($Y$) que tienen:\par\smallskip{\centering {\small$\begin{array}{|c|c|c|c|}\hline X\backslash Y & 0 & 1 & 2\\\hline 1 & 4 & 2 & 0\\ 2 & 3 & 6 & 3\\ 3 & 0 & 3 & 4\\\hline \end{array}$}\par}\smallskip\begin{tasks}(1)
  \task Añade las distribuciones marginales a la tabla.
  \task ¿Cuántas familias tienen 2 hijos y 1 coche? ¿Cuántas tienen más de uno?
  \task Calcula las medias marginales $\bar x$ y $\bar y$.
\end{tasks}\recorda{En una tabla de doble entrada, las sumas de las filas dan la distribución marginal de $X$ y las de las columnas, la de $Y$. La suma de todas las casillas es $N$.}

\ex{2} Las notas de Matemáticas ($X$) y de Física ($Y$) de 12 alumnos, en una escala de 1 a 3 (1 = insuficiente, 2 = aprobado, 3 = notable o más), han sido: \ $(2,1)$, $(3,2)$, $(2,2)$, $(1,1)$, $(3,3)$, $(2,1)$, $(1,1)$, $(2,2)$, $(3,2)$, $(2,2)$, $(1,2)$, $(3,3)$. \ Construye la tabla de doble entrada y calcula las medias marginales.

\ex{2} Las horas de estudio semanales ($X$) y la nota final ($Y$) de 50 alumnos, agrupadas en intervalos, se recogen en esta tabla:\par\smallskip{\centering {\small$\begin{array}{|c|c|c|c|}\hline X\backslash Y & {[0,5)} & {[5,7)} & {[7,10]}\\\hline {[0,2)} & 8 & 4 & 1\\ {[2,4)} & 5 & 10 & 6\\ {[4,6)} & 1 & 5 & 10\\\hline \end{array}$}\par}\smallskip\begin{tasks}(1)
  \task Añade las distribuciones marginales y calcula las medias marginales (usa las marcas de clase).
  \task ¿Qué porcentaje de alumnos estudia entre 2 y 4 horas y saca un 7 o más?
  \task De los alumnos que estudian 4 horas o más, ¿qué porcentaje saca un 7 o más? ¿Y de los que estudian menos de 2 horas?
  \task De los alumnos que han suspendido (nota menor que 5), ¿qué porcentaje estudia menos de 2 horas?
\end{tasks}\recorda{La marca de clase es el punto medio del intervalo. Una distribución condicionada es una fila o una columna: los porcentajes se calculan respecto del total de esa fila o columna.}


% =====================================================
\seccio{2}{Nube de puntos}

\exf{1} Un coche gasta gasolina y anotamos el combustible que queda en el depósito ($Y$, en litros) al final de cada hora de viaje ($X$): \par\smallskip{\centering {\small$\begin{array}{|c|cccccccc|}\hline x_i & 1 & 2 & 3 & 4 & 5 & 6 & 7 & 8\\\hline y_i & 62 & 60 & 55 & 53 & 50 & 46 & 44 & 40\\\hline\end{array}$}\par}\smallskip Representa la nube de puntos y el centro de gravedad. ¿Qué tipo de correlación se observa?\recorda{Cada par $(x_i,y_i)$ es un punto. El centro de gravedad es $(\bar x,\bar y)$.}

\ex{1} Asocia cada nube con uno de estos coeficientes de correlación: \ $r=0$, \ $r=0{,}9$, \ $r=-0{,}5$, \ $r=-0{,}95$.\begin{tasks}(4)
  \task \par\begin{tikzpicture}[scale=0.22,baseline=(current bounding box.center)]\draw[muted!30,very thin] (0,0) grid (10.0,10.0);\draw[-latex,thick] (0,0)--(10.5,0) node[right,font=\scriptsize]{$X$};\draw[-latex,thick] (0,0)--(0,10.5) node[above,font=\scriptsize]{$Y$};\node[below,font=\tiny] at (0.0,0) {$0$};\node[below,font=\tiny] at (2.0,0) {$2$};\node[below,font=\tiny] at (4.0,0) {$4$};\node[below,font=\tiny] at (6.0,0) {$6$};\node[below,font=\tiny] at (8.0,0) {$8$};\node[below,font=\tiny] at (10.0,0) {$10$};\node[left,font=\tiny] at (0,0.0) {$0$};\node[left,font=\tiny] at (0,2.0) {$2$};\node[left,font=\tiny] at (0,4.0) {$4$};\node[left,font=\tiny] at (0,6.0) {$6$};\node[left,font=\tiny] at (0,8.0) {$8$};\node[left,font=\tiny] at (0,10.0) {$10$};\fill[accent] (6.968,7.394) circle (0.16);\fill[accent] (7.235,7.358) circle (0.16);\fill[accent] (4.95,4.95) circle (0.16);\fill[accent] (3.577,2.806) circle (0.16);\fill[accent] (3.036,3.567) circle (0.16);\fill[accent] (4.892,4.859) circle (0.16);\fill[accent] (3.152,3.886) circle (0.16);\fill[accent] (2.467,2.897) circle (0.16);\fill[accent] (5.169,5.907) circle (0.16);\fill[accent] (5.061,4.9) circle (0.16);\fill[accent] (5.743,5.687) circle (0.16);\fill[accent] (3.33,3.996) circle (0.16);\fill[accent] (4.848,3.889) circle (0.16);\fill[accent] (4.733,4.34) circle (0.16);\fill[accent] (2.353,2.225) circle (0.16);\fill[accent] (5.729,7.1) circle (0.16);\fill[accent] (5.37,5.132) circle (0.16);\fill[accent] (8.787,8.651) circle (0.16);\fill[accent] (5.175,5.536) circle (0.16);\fill[accent] (4.601,4.324) circle (0.16);\fill[accent] (6.877,5.252) circle (0.16);\fill[accent] (5.169,5.807) circle (0.16);\fill[accent] (6.342,5.711) circle (0.16);\fill[accent] (4.236,4.812) circle (0.16);\fill[accent] (5.201,4.015) circle (0.16);\end{tikzpicture}
  \task \par\begin{tikzpicture}[scale=0.22,baseline=(current bounding box.center)]\draw[muted!30,very thin] (0,0) grid (10.0,10.0);\draw[-latex,thick] (0,0)--(10.5,0) node[right,font=\scriptsize]{$X$};\draw[-latex,thick] (0,0)--(0,10.5) node[above,font=\scriptsize]{$Y$};\node[below,font=\tiny] at (0.0,0) {$0$};\node[below,font=\tiny] at (2.0,0) {$2$};\node[below,font=\tiny] at (4.0,0) {$4$};\node[below,font=\tiny] at (6.0,0) {$6$};\node[below,font=\tiny] at (8.0,0) {$8$};\node[below,font=\tiny] at (10.0,0) {$10$};\node[left,font=\tiny] at (0,0.0) {$0$};\node[left,font=\tiny] at (0,2.0) {$2$};\node[left,font=\tiny] at (0,4.0) {$4$};\node[left,font=\tiny] at (0,6.0) {$6$};\node[left,font=\tiny] at (0,8.0) {$8$};\node[left,font=\tiny] at (0,10.0) {$10$};\fill[accent] (8.913,1.832) circle (0.16);\fill[accent] (4.499,5.153) circle (0.16);\fill[accent] (6.055,4.023) circle (0.16);\fill[accent] (5.69,4.471) circle (0.16);\fill[accent] (6.702,3.546) circle (0.16);\fill[accent] (3.412,6.621) circle (0.16);\fill[accent] (4.864,5.466) circle (0.16);\fill[accent] (4.369,4.369) circle (0.16);\fill[accent] (3.894,6.15) circle (0.16);\fill[accent] (4.233,5.793) circle (0.16);\fill[accent] (4.72,5.205) circle (0.16);\fill[accent] (5.052,5.634) circle (0.16);\fill[accent] (4.141,5.569) circle (0.16);\fill[accent] (6.095,3.981) circle (0.16);\fill[accent] (4.669,4.638) circle (0.16);\fill[accent] (0.771,9.384) circle (0.16);\fill[accent] (7.225,2.267) circle (0.16);\fill[accent] (4.898,4.589) circle (0.16);\fill[accent] (4.381,6.62) circle (0.16);\fill[accent] (5.869,4.509) circle (0.16);\fill[accent] (5.812,4.288) circle (0.16);\fill[accent] (5.552,4.616) circle (0.16);\fill[accent] (4.217,5.309) circle (0.16);\fill[accent] (5.756,3.938) circle (0.16);\fill[accent] (3.213,7.027) circle (0.16);\end{tikzpicture}
  \task \par\begin{tikzpicture}[scale=0.22,baseline=(current bounding box.center)]\draw[muted!30,very thin] (0,0) grid (10.0,10.0);\draw[-latex,thick] (0,0)--(10.5,0) node[right,font=\scriptsize]{$X$};\draw[-latex,thick] (0,0)--(0,10.5) node[above,font=\scriptsize]{$Y$};\node[below,font=\tiny] at (0.0,0) {$0$};\node[below,font=\tiny] at (2.0,0) {$2$};\node[below,font=\tiny] at (4.0,0) {$4$};\node[below,font=\tiny] at (6.0,0) {$6$};\node[below,font=\tiny] at (8.0,0) {$8$};\node[below,font=\tiny] at (10.0,0) {$10$};\node[left,font=\tiny] at (0,0.0) {$0$};\node[left,font=\tiny] at (0,2.0) {$2$};\node[left,font=\tiny] at (0,4.0) {$4$};\node[left,font=\tiny] at (0,6.0) {$6$};\node[left,font=\tiny] at (0,8.0) {$8$};\node[left,font=\tiny] at (0,10.0) {$10$};\fill[accent] (5.371,5.94) circle (0.16);\fill[accent] (7.362,6.727) circle (0.16);\fill[accent] (3.602,4.307) circle (0.16);\fill[accent] (6.918,5.012) circle (0.16);\fill[accent] (4.761,2.585) circle (0.16);\fill[accent] (4.757,4.83) circle (0.16);\fill[accent] (8.482,2.363) circle (0.16);\fill[accent] (5.479,3.479) circle (0.16);\fill[accent] (5.134,7.219) circle (0.16);\fill[accent] (6.465,2.33) circle (0.16);\fill[accent] (7.15,6.795) circle (0.16);\fill[accent] (5.154,6.07) circle (0.16);\fill[accent] (6.221,5.133) circle (0.16);\fill[accent] (3.529,6.245) circle (0.16);\fill[accent] (4.575,6.359) circle (0.16);\fill[accent] (4.452,7.127) circle (0.16);\fill[accent] (2.911,5.213) circle (0.16);\fill[accent] (2.608,4.671) circle (0.16);\fill[accent] (2.404,4.65) circle (0.16);\fill[accent] (4.796,4.114) circle (0.16);\fill[accent] (4.91,5.513) circle (0.16);\fill[accent] (4.655,4.41) circle (0.16);\fill[accent] (5.327,7.173) circle (0.16);\fill[accent] (2.906,2.778) circle (0.16);\fill[accent] (5.071,3.958) circle (0.16);\end{tikzpicture}
  \task \par\begin{tikzpicture}[scale=0.22,baseline=(current bounding box.center)]\draw[muted!30,very thin] (0,0) grid (10.0,10.0);\draw[-latex,thick] (0,0)--(10.5,0) node[right,font=\scriptsize]{$X$};\draw[-latex,thick] (0,0)--(0,10.5) node[above,font=\scriptsize]{$Y$};\node[below,font=\tiny] at (0.0,0) {$0$};\node[below,font=\tiny] at (2.0,0) {$2$};\node[below,font=\tiny] at (4.0,0) {$4$};\node[below,font=\tiny] at (6.0,0) {$6$};\node[below,font=\tiny] at (8.0,0) {$8$};\node[below,font=\tiny] at (10.0,0) {$10$};\node[left,font=\tiny] at (0,0.0) {$0$};\node[left,font=\tiny] at (0,2.0) {$2$};\node[left,font=\tiny] at (0,4.0) {$4$};\node[left,font=\tiny] at (0,6.0) {$6$};\node[left,font=\tiny] at (0,8.0) {$8$};\node[left,font=\tiny] at (0,10.0) {$10$};\fill[accent] (4.938,4.719) circle (0.16);\fill[accent] (5.633,3.133) circle (0.16);\fill[accent] (4.115,3.9) circle (0.16);\fill[accent] (5.449,4.061) circle (0.16);\fill[accent] (6.39,4.807) circle (0.16);\fill[accent] (5.546,4.358) circle (0.16);\fill[accent] (7.433,4.713) circle (0.16);\fill[accent] (3.42,8.505) circle (0.16);\fill[accent] (4.982,3.89) circle (0.16);\fill[accent] (3.715,4.178) circle (0.16);\fill[accent] (3.586,5.108) circle (0.16);\fill[accent] (4.57,7.04) circle (0.16);\fill[accent] (5.234,3.982) circle (0.16);\fill[accent] (5.558,7.351) circle (0.16);\fill[accent] (5.704,2.901) circle (0.16);\fill[accent] (8.534,2.675) circle (0.16);\fill[accent] (6.286,4.782) circle (0.16);\fill[accent] (2.256,4.433) circle (0.16);\fill[accent] (5.206,4.13) circle (0.16);\fill[accent] (3.849,6.173) circle (0.16);\fill[accent] (4.025,6.571) circle (0.16);\fill[accent] (7.013,4.881) circle (0.16);\fill[accent] (4.505,4.849) circle (0.16);\fill[accent] (1.671,8.127) circle (0.16);\fill[accent] (5.38,5.732) circle (0.16);\end{tikzpicture}
\end{tasks}

\ex{1} Indica, para cada par de variables, si esperas una correlación positiva, negativa o ninguna correlación:\begin{tasks}(1)
  \task La altura y el peso de 30 adolescentes.
  \task El precio de una entrada de cine y el número de espectadores.
  \task La nota de inglés y el número de zapato.
  \task Los kilómetros recorridos y la gasolina gastada.
  \task La temperatura exterior y el consumo de calefacción.
  \task La renta de una familia y su gasto en viajes.
\end{tasks}


% =====================================================
\seccio{3}{Correlación y regresión. Tipos de dependencia}

\exf{1} Di si la dependencia entre las variables es funcional, estadística o si son independientes:\begin{tasks}(1)
  \task El lado de un cuadrado y su perímetro.
  \task Las horas de sol de un día y las ventas de crema solar.
  \task El número de DNI y la altura.
  \task Los años de experiencia y el sueldo de un trabajador.
  \task El precio de un trayecto en taxi y los kilómetros (tarifa fija).
  \task Las horas de entrenamiento y el tiempo en una carrera.
\end{tasks}\recorda{\textbf{Funcional}: una variable determina exactamente la otra ($y=f(x)$). \textbf{Estadística}: hay relación, pero no exacta (nube alrededor de una línea). \textbf{Independientes}: no hay ninguna relación.}

\ex{1} ¿Qué tipo de dependencia muestra cada nube de puntos? ¿En cuál tiene sentido buscar una recta de regresión para hacer estimaciones?\begin{tasks}(3)
  \task \par\begin{tikzpicture}[scale=0.22,baseline=(current bounding box.center)]\draw[muted!30,very thin] (0,0) grid (10.0,10.0);\draw[-latex,thick] (0,0)--(10.5,0) node[right,font=\scriptsize]{$X$};\draw[-latex,thick] (0,0)--(0,10.5) node[above,font=\scriptsize]{$Y$};\node[below,font=\tiny] at (0.0,0) {$0$};\node[below,font=\tiny] at (2.0,0) {$2$};\node[below,font=\tiny] at (4.0,0) {$4$};\node[below,font=\tiny] at (6.0,0) {$6$};\node[below,font=\tiny] at (8.0,0) {$8$};\node[below,font=\tiny] at (10.0,0) {$10$};\node[left,font=\tiny] at (0,0.0) {$0$};\node[left,font=\tiny] at (0,2.0) {$2$};\node[left,font=\tiny] at (0,4.0) {$4$};\node[left,font=\tiny] at (0,6.0) {$6$};\node[left,font=\tiny] at (0,8.0) {$8$};\node[left,font=\tiny] at (0,10.0) {$10$};\fill[accent] (0.5,0.525) circle (0.16);\fill[accent] (0.95,0.59) circle (0.16);\fill[accent] (1.4,0.696) circle (0.16);\fill[accent] (1.85,0.842) circle (0.16);\fill[accent] (2.3,1.029) circle (0.16);\fill[accent] (2.75,1.256) circle (0.16);\fill[accent] (3.2,1.524) circle (0.16);\fill[accent] (3.65,1.832) circle (0.16);\fill[accent] (4.1,2.181) circle (0.16);\fill[accent] (4.55,2.57) circle (0.16);\fill[accent] (5.0,3.0) circle (0.16);\fill[accent] (5.45,3.47) circle (0.16);\fill[accent] (5.9,3.981) circle (0.16);\fill[accent] (6.35,4.532) circle (0.16);\fill[accent] (6.8,5.124) circle (0.16);\fill[accent] (7.25,5.756) circle (0.16);\fill[accent] (7.7,6.429) circle (0.16);\fill[accent] (8.15,7.142) circle (0.16);\fill[accent] (8.6,7.896) circle (0.16);\fill[accent] (9.05,8.69) circle (0.16);\fill[accent] (9.5,9.525) circle (0.16);\end{tikzpicture}
  \task \par\begin{tikzpicture}[scale=0.22,baseline=(current bounding box.center)]\draw[muted!30,very thin] (0,0) grid (10.0,10.0);\draw[-latex,thick] (0,0)--(10.5,0) node[right,font=\scriptsize]{$X$};\draw[-latex,thick] (0,0)--(0,10.5) node[above,font=\scriptsize]{$Y$};\node[below,font=\tiny] at (0.0,0) {$0$};\node[below,font=\tiny] at (2.0,0) {$2$};\node[below,font=\tiny] at (4.0,0) {$4$};\node[below,font=\tiny] at (6.0,0) {$6$};\node[below,font=\tiny] at (8.0,0) {$8$};\node[below,font=\tiny] at (10.0,0) {$10$};\node[left,font=\tiny] at (0,0.0) {$0$};\node[left,font=\tiny] at (0,2.0) {$2$};\node[left,font=\tiny] at (0,4.0) {$4$};\node[left,font=\tiny] at (0,6.0) {$6$};\node[left,font=\tiny] at (0,8.0) {$8$};\node[left,font=\tiny] at (0,10.0) {$10$};\fill[accent] (5.84,6.51) circle (0.16);\fill[accent] (2.479,3.204) circle (0.16);\fill[accent] (3.968,3.37) circle (0.16);\fill[accent] (5.214,5.642) circle (0.16);\fill[accent] (7.272,7.225) circle (0.16);\fill[accent] (5.113,4.699) circle (0.16);\fill[accent] (2.67,2.593) circle (0.16);\fill[accent] (5.567,6.499) circle (0.16);\fill[accent] (3.347,3.333) circle (0.16);\fill[accent] (6.907,6.978) circle (0.16);\fill[accent] (4.731,5.891) circle (0.16);\fill[accent] (7.722,6.158) circle (0.16);\fill[accent] (4.915,6.354) circle (0.16);\fill[accent] (3.601,4.649) circle (0.16);\fill[accent] (2.948,2.807) circle (0.16);\fill[accent] (4.577,3.936) circle (0.16);\fill[accent] (5.76,4.719) circle (0.16);\fill[accent] (6.853,7.316) circle (0.16);\fill[accent] (5.498,3.445) circle (0.16);\fill[accent] (4.022,4.938) circle (0.16);\fill[accent] (5.842,5.87) circle (0.16);\fill[accent] (2.997,2.819) circle (0.16);\fill[accent] (5.848,5.347) circle (0.16);\fill[accent] (4.005,3.912) circle (0.16);\fill[accent] (7.303,6.785) circle (0.16);\end{tikzpicture}
  \task \par\begin{tikzpicture}[scale=0.22,baseline=(current bounding box.center)]\draw[muted!30,very thin] (0,0) grid (10.0,10.0);\draw[-latex,thick] (0,0)--(10.5,0) node[right,font=\scriptsize]{$X$};\draw[-latex,thick] (0,0)--(0,10.5) node[above,font=\scriptsize]{$Y$};\node[below,font=\tiny] at (0.0,0) {$0$};\node[below,font=\tiny] at (2.0,0) {$2$};\node[below,font=\tiny] at (4.0,0) {$4$};\node[below,font=\tiny] at (6.0,0) {$6$};\node[below,font=\tiny] at (8.0,0) {$8$};\node[below,font=\tiny] at (10.0,0) {$10$};\node[left,font=\tiny] at (0,0.0) {$0$};\node[left,font=\tiny] at (0,2.0) {$2$};\node[left,font=\tiny] at (0,4.0) {$4$};\node[left,font=\tiny] at (0,6.0) {$6$};\node[left,font=\tiny] at (0,8.0) {$8$};\node[left,font=\tiny] at (0,10.0) {$10$};\fill[accent] (4.681,5.786) circle (0.16);\fill[accent] (6.132,4.091) circle (0.16);\fill[accent] (4.737,8.063) circle (0.16);\fill[accent] (4.569,6.056) circle (0.16);\fill[accent] (3.407,7.061) circle (0.16);\fill[accent] (4.762,4.05) circle (0.16);\fill[accent] (7.267,4.049) circle (0.16);\fill[accent] (5.967,4.623) circle (0.16);\fill[accent] (7.125,5.125) circle (0.16);\fill[accent] (5.635,6.272) circle (0.16);\fill[accent] (5.911,5.636) circle (0.16);\fill[accent] (5.515,4.408) circle (0.16);\fill[accent] (2.015,3.248) circle (0.16);\fill[accent] (6.781,4.386) circle (0.16);\fill[accent] (6.122,7.323) circle (0.16);\fill[accent] (6.108,3.837) circle (0.16);\fill[accent] (1.967,5.308) circle (0.16);\fill[accent] (1.868,5.612) circle (0.16);\fill[accent] (3.483,2.452) circle (0.16);\fill[accent] (4.28,5.159) circle (0.16);\fill[accent] (5.742,7.439) circle (0.16);\fill[accent] (5.078,1.681) circle (0.16);\fill[accent] (6.15,4.676) circle (0.16);\fill[accent] (3.951,4.867) circle (0.16);\fill[accent] (5.748,3.792) circle (0.16);\end{tikzpicture}
\end{tasks}

\ex{2} En un estudio de varias ciudades se ha encontrado una correlación positiva fuerte entre el número de bomberos que actúan en un incendio y los daños que causa. ¿Quiere decir que los bomberos causan los daños? ¿Qué explicación puede haber?


% =====================================================
\seccio{4}{Correlación lineal. Covarianza}

\exf{2} Para la variable bidimensional siguiente, haz la tabla de cálculos y halla las medias, las varianzas, las desviaciones típicas y la covarianza:\par\smallskip{\centering {\small$\begin{array}{|c|ccccc|}\hline x_i & 2 & 4 & 5 & 6 & 8\\\hline y_i & 3 & 5 & 4 & 7 & 9\\\hline\end{array}$}\par}\recorda{$\sigma_x^2=\frac{\sum x_i^2}{N}-\bar x^2$, \ $\sigma_y^2=\frac{\sum y_i^2}{N}-\bar y^2$, \ $\sigma_{xy}=\frac{\sum x_iy_i}{N}-\bar x\,\bar y$.}

\ex{2} Calcula la covarianza de los datos siguientes e interpreta su signo:\par\smallskip{\centering {\small$\begin{array}{|c|ccccc|}\hline x_i & 10 & 20 & 30 & 40 & 50\\\hline y_i & 9 & 8 & 6 & 5 & 2\\\hline\end{array}$}\par}

\ex{2} Calcula la covarianza de las notas de Matemáticas y de Física del ejercicio 2 (la tabla de doble entrada de 12 alumnos).

\ex{2} Calcula la covarianza de la variable bidimensional del ejercicio 1 (hijos y coches).

\ex{3} Calcula las varianzas y la covarianza de las horas de estudio y las notas del ejercicio 3 (datos agrupados). ¿Cómo es la relación?\recorda{Con datos agrupados, cada casilla aporta $x_i\,y_j\,f_{ij}$ (con las marcas de clase): $\sigma_{xy}=\frac{\sum x_iy_jf_{ij}}{N}-\bar x\,\bar y$.}


% =====================================================
\seccio{5}{Coeficiente de correlación lineal}

\exf{2} Calcula el coeficiente de correlación de los datos de los ejercicios de la sección 4 e interprétalo.\begin{tasks}(1)
  \task Los datos del primer ejercicio de la sección 4.
  \task Los datos del segundo ejercicio de la sección 4.
  \task Hijos y coches (ejercicio 1).
  \task Horas de estudio y notas agrupadas (ejercicio 3).
\end{tasks}\recorda{$r=\dfrac{\sigma_{xy}}{\sigma_x\,\sigma_y}$, \ $-1\leq r\leq 1$. El signo da el sentido de la relación y $|r|$, la fuerza: cerca de 1, fuerte; cerca de 0, no hay correlación lineal.}

\ex{2} La altura (cm) y el peso (kg) de 7 personas han sido:\par\smallskip{\centering {\small$\begin{array}{|c|ccccccc|}\hline x_i\ (\text{cm}) & 150 & 155 & 160 & 165 & 170 & 175 & 180\\\hline y_i\ (\text{kg}) & 48 & 52 & 55 & 61 & 63 & 70 & 72\\\hline\end{array}$}\par}\smallskip Representa la nube de puntos, calcula el coeficiente de correlación e interprétalo.

\ex{1} Si en el ejercicio anterior expresamos la altura en metros en lugar de centímetros, ¿qué les pasa a la covarianza y al coeficiente de correlación?

\ex{2} Las horas semanales de videojuegos ($X$) y la nota de un examen ($Y$) de 6 alumnos han sido:\par\smallskip{\centering {\small$\begin{array}{|c|cccccc|}\hline x_i & 0 & 1 & 2 & 3 & 4 & 5\\\hline y_i & 10 & 8 & 9 & 6 & 4 & 3\\\hline\end{array}$}\par}\smallskip Calcula el coeficiente de correlación. ¿Podemos afirmar que jugar a videojuegos hace bajar la nota?

\ex{2} Una heladería anota la temperatura máxima ($X$, en °C) y los helados vendidos ($Y$) durante seis días:\par\smallskip{\centering {\small$\begin{array}{|c|cccccc|}\hline x_i & 18 & 21 & 24 & 27 & 30 & 33\\\hline y_i & 35 & 52 & 58 & 75 & 84 & 102\\\hline\end{array}$}\par}\smallskip Calcula el coeficiente de correlación e interprétalo.

\ex{2} Calcula el coeficiente de correlación entre las notas de Matemáticas y de Física del ejercicio 2.

\ex{1} Si el coeficiente de correlación entre dos variables es $r=0{,}8$ y las desviaciones típicas son $\sigma_x=2$ y $\sigma_y=5$, ¿cuánto vale la covarianza?

\ex{1} Si $\sigma_{xy}=-12$, $\sigma_x=4$ y $\sigma_y=5$, ¿cuánto vale $r$? ¿Y si $\sigma_{xy}=3$, $\sigma_x=2$ y $\sigma_y=1$? ¿Es posible?

\ex{3} Una variable bidimensional tiene $N=10$, $\sum x_i=50$, $\sum y_i=80$, $\sum x_i^2=290$, $\sum y_i^2=730$ y $\sum x_iy_i=440$. Calcula $r$.

\ex{2} Observa estas nubes de puntos. Describe la relación que hay entre las variables y di si el coeficiente de correlación lineal será próximo a 0 o no.\begin{tasks}(3)
  \task \par\begin{tikzpicture}[scale=0.22,baseline=(current bounding box.center)]\draw[muted!30,very thin] (0,0) grid (10.0,10.0);\draw[-latex,thick] (0,0)--(10.5,0) node[right,font=\scriptsize]{$X$};\draw[-latex,thick] (0,0)--(0,10.5) node[above,font=\scriptsize]{$Y$};\node[below,font=\tiny] at (0.0,0) {$0$};\node[below,font=\tiny] at (2.0,0) {$2$};\node[below,font=\tiny] at (4.0,0) {$4$};\node[below,font=\tiny] at (6.0,0) {$6$};\node[below,font=\tiny] at (8.0,0) {$8$};\node[below,font=\tiny] at (10.0,0) {$10$};\node[left,font=\tiny] at (0,0.0) {$0$};\node[left,font=\tiny] at (0,2.0) {$2$};\node[left,font=\tiny] at (0,4.0) {$4$};\node[left,font=\tiny] at (0,6.0) {$6$};\node[left,font=\tiny] at (0,8.0) {$8$};\node[left,font=\tiny] at (0,10.0) {$10$};\fill[accent] (0.3,1.416) circle (0.16);\fill[accent] (0.7,2.819) circle (0.16);\fill[accent] (1.1,4.031) circle (0.16);\fill[accent] (1.5,5.243) circle (0.16);\fill[accent] (1.9,5.924) circle (0.16);\fill[accent] (2.3,6.955) circle (0.16);\fill[accent] (2.7,6.583) circle (0.16);\fill[accent] (3.1,7.695) circle (0.16);\fill[accent] (3.5,8.745) circle (0.16);\fill[accent] (3.9,8.755) circle (0.16);\fill[accent] (4.3,9.31) circle (0.16);\fill[accent] (4.7,8.504) circle (0.16);\fill[accent] (5.1,8.959) circle (0.16);\fill[accent] (5.5,8.608) circle (0.16);\fill[accent] (5.9,8.769) circle (0.16);\fill[accent] (6.3,8.497) circle (0.16);\fill[accent] (6.7,7.404) circle (0.16);\fill[accent] (7.1,7.117) circle (0.16);\fill[accent] (7.5,6.548) circle (0.16);\fill[accent] (7.9,6.556) circle (0.16);\fill[accent] (8.3,5.507) circle (0.16);\fill[accent] (8.7,3.8) circle (0.16);\fill[accent] (9.1,3.473) circle (0.16);\fill[accent] (9.5,1.479) circle (0.16);\end{tikzpicture}
  \task \par\begin{tikzpicture}[scale=0.22,baseline=(current bounding box.center)]\draw[muted!30,very thin] (0,0) grid (10.0,10.0);\draw[-latex,thick] (0,0)--(10.5,0) node[right,font=\scriptsize]{$X$};\draw[-latex,thick] (0,0)--(0,10.5) node[above,font=\scriptsize]{$Y$};\node[below,font=\tiny] at (0.0,0) {$0$};\node[below,font=\tiny] at (2.0,0) {$2$};\node[below,font=\tiny] at (4.0,0) {$4$};\node[below,font=\tiny] at (6.0,0) {$6$};\node[below,font=\tiny] at (8.0,0) {$8$};\node[below,font=\tiny] at (10.0,0) {$10$};\node[left,font=\tiny] at (0,0.0) {$0$};\node[left,font=\tiny] at (0,2.0) {$2$};\node[left,font=\tiny] at (0,4.0) {$4$};\node[left,font=\tiny] at (0,6.0) {$6$};\node[left,font=\tiny] at (0,8.0) {$8$};\node[left,font=\tiny] at (0,10.0) {$10$};\fill[accent] (0.3,0.229) circle (0.16);\fill[accent] (0.7,0.015) circle (0.16);\fill[accent] (1.5,0.467) circle (0.16);\fill[accent] (1.9,0.188) circle (0.16);\fill[accent] (2.3,0.252) circle (0.16);\fill[accent] (2.7,0.707) circle (0.16);\fill[accent] (3.1,0.738) circle (0.16);\fill[accent] (3.5,0.547) circle (0.16);\fill[accent] (3.9,1.002) circle (0.16);\fill[accent] (4.3,0.933) circle (0.16);\fill[accent] (4.7,1.169) circle (0.16);\fill[accent] (5.1,1.13) circle (0.16);\fill[accent] (5.5,1.732) circle (0.16);\fill[accent] (5.9,1.883) circle (0.16);\fill[accent] (6.3,2.348) circle (0.16);\fill[accent] (6.7,2.698) circle (0.16);\fill[accent] (7.1,2.858) circle (0.16);\fill[accent] (7.5,3.431) circle (0.16);\fill[accent] (7.9,3.974) circle (0.16);\fill[accent] (8.3,4.721) circle (0.16);\fill[accent] (8.7,5.577) circle (0.16);\fill[accent] (9.1,6.755) circle (0.16);\fill[accent] (9.5,8.16) circle (0.16);\end{tikzpicture}
  \task \par\begin{tikzpicture}[scale=0.22,baseline=(current bounding box.center)]\draw[muted!30,very thin] (0,0) grid (10.0,10.0);\draw[-latex,thick] (0,0)--(10.5,0) node[right,font=\scriptsize]{$X$};\draw[-latex,thick] (0,0)--(0,10.5) node[above,font=\scriptsize]{$Y$};\node[below,font=\tiny] at (0.0,0) {$0$};\node[below,font=\tiny] at (2.0,0) {$2$};\node[below,font=\tiny] at (4.0,0) {$4$};\node[below,font=\tiny] at (6.0,0) {$6$};\node[below,font=\tiny] at (8.0,0) {$8$};\node[below,font=\tiny] at (10.0,0) {$10$};\node[left,font=\tiny] at (0,0.0) {$0$};\node[left,font=\tiny] at (0,2.0) {$2$};\node[left,font=\tiny] at (0,4.0) {$4$};\node[left,font=\tiny] at (0,6.0) {$6$};\node[left,font=\tiny] at (0,8.0) {$8$};\node[left,font=\tiny] at (0,10.0) {$10$};\fill[accent] (1.0,0.604) circle (0.16);\fill[accent] (1.15,1.414) circle (0.16);\fill[accent] (1.3,1.005) circle (0.16);\fill[accent] (1.45,0.972) circle (0.16);\fill[accent] (1.6,1.582) circle (0.16);\fill[accent] (1.75,1.177) circle (0.16);\fill[accent] (1.9,0.979) circle (0.16);\fill[accent] (2.05,1.177) circle (0.16);\fill[accent] (9.0,9.0) circle (0.16);\end{tikzpicture}
\end{tasks}

\ex{1} Razona si son verdaderas o falsas: \ a) Si $r=0$, las variables no tienen ninguna relación. \ b) Si $r=-0{,}98$, la relación es muy débil porque $r$ es negativo. \ c) Si todos los puntos están sobre una recta creciente, $r=1$.


% =====================================================
\seccio{6}{Rectas de regresión}

\exf{2} Con los datos del primer ejercicio de la sección 4, halla la recta de regresión de $Y$ sobre $X$ y la de $X$ sobre $Y$. Estima el valor de $y$ para $x=7$ y el de $x$ para $y=6$.\recorda{Recta de regresión de $Y$ sobre $X$: \ $y-\bar y=\dfrac{\sigma_{xy}}{\sigma_x^2}(x-\bar x)$ (para estimar $y$). \ De $X$ sobre $Y$: \ $x-\bar x=\dfrac{\sigma_{xy}}{\sigma_y^2}(y-\bar y)$ (para estimar $x$). Las estimaciones son fiables si $|r|$ está cerca de 1 y no nos salimos del intervalo de los datos.}

\ex{2} Una empresa ha anotado durante 8 meses la inversión en publicidad ($X$, en miles de euros) y las ventas ($Y$, en miles de euros):\par\smallskip{\centering {\small$\begin{array}{|c|cccccccc|}\hline x_i & 1 & 2 & 3 & 4 & 5 & 6 & 7 & 8\\\hline y_i & 12 & 15 & 15 & 19 & 22 & 22 & 26 & 27\\\hline\end{array}$}\par}\smallskip\begin{tasks}(1)
  \task Calcula el coeficiente de correlación e interprétalo.
  \task Halla la recta de regresión de $Y$ sobre $X$. ¿Qué significa su pendiente?
  \task ¿Qué ventas se esperan si se invierten 5500 € en publicidad?
  \task ¿Cuánto hay que invertir en publicidad para vender 24\,000 €? (Usa la recta adecuada.)
\end{tasks}

\ex{2} Con las alturas y los pesos de la sección 5, halla la recta de regresión del peso sobre la altura y la de la altura sobre el peso. ¿Qué peso se estima para una persona de 168 cm? ¿Y para un niño de 120 cm? ¿Qué altura se estima para una persona de 65 kg?

\ex{2} Con los datos de la heladería de la sección 5, ¿cuántos helados se estima que se venderán un día con una máxima de 25 °C?

\ex{2} Con la tabla de hijos y coches del ejercicio 1, halla la recta de regresión de $Y$ sobre $X$ y estima el número de coches de una familia con 2 hijos. ¿Es fiable?

\ex{2} Con las sumas del ejercicio de la sección 5 ($N=10$, $\sum x_i=50$, $\sum y_i=80$, $\sum x_i^2=290$, $\sum x_iy_i=440$), halla la recta de regresión de $Y$ sobre $X$ y estima $y$ para $x=6$.

\ex{2} De una distribución bidimensional sabemos que $\bar x=8$, $\bar y=15$, $\sigma_x=2$, $\sigma_y=4$ y $r=-0{,}75$. Halla las dos rectas de regresión y estima $y$ para $x=9$ y $x$ para $y=12$.

\ex{3} Las rectas de regresión de una variable bidimensional son $y=0{,}5x+2$ (de $Y$ sobre $X$) y $x=1{,}2y-1$ (de $X$ sobre $Y$). Calcula $\bar x$, $\bar y$ y el coeficiente de correlación.\recorda{Las dos rectas pasan por $(\bar x,\bar y)$ y el producto de los coeficientes de regresión es $m_{yx}\cdot m_{xy}=\frac{\sigma_{xy}}{\sigma_x^2}\cdot\frac{\sigma_{xy}}{\sigma_y^2}=r^2$.}

\ex{3} Las rectas de regresión de una distribución son $4x-5y=0$ y $2x-y=6$. Halla el centro de gravedad, averigua qué recta es la de $Y$ sobre $X$ y calcula $r$. Estima $y$ para $x=6$ y $x$ para $y=5$.

\ex{2} Razona si pueden ser las rectas de regresión de una distribución: \ a) $y=2x+1$ e $y=-x+4$; \ b) $y=-0{,}4x+6$ (de $Y$ sobre $X$) y $x=-0{,}9y+10$ (de $X$ sobre $Y$).

\vfill
\begin{center}\color{muted}\small\faCheckCircle[regular]\ \ ¡Ánimo! Comprueba tus resultados con la calculadora bidimensional de los apuntes interactivos.\end{center}

\end{document}
